\documentclass{beamer}
% Theme selection
\usetheme{Madrid}
\usecolortheme{seagull}
\setbeamertemplate{footline}[frame number]{}
\setbeamertemplate{navigation symbols}{}

% Packages
\usepackage[utf8]{inputenc}
\usepackage{graphicx}
\usepackage{amsmath}
\usepackage{hyperref}
\usepackage[T1]{fontenc}

% Presentation metadata$
\title{Aproksymacja średniokwadratowa ciągła w przestrzeni $L_w^2(-1,1)$ dla $w(x) = 1$ w bazie wielomianów Legendre'a}
\subtitle{}
\author{Karol Sęk\\~\\
334831}
\institute{grupa 5c\\~\\
czwartek 13:15}
\date{projekt 1\\~\\
zadanie 51}
\setbeamerfont{institute}{size=\normalsize}
\setbeamerfont{date}{size=\normalsize}
\setbeamerfont{author}{size=\normalsize}


\begin{document}

% Title slide
\frame{\titlepage}

% Table of contents
\begin{frame}
\frametitle{Spis treści}
\tableofcontents
\end{frame}

% Section 1
\section{Wprowadzenie}
\begin{frame}
\frametitle{Wprowadzenie}
\begin{alertblock}{Problem}
    Niech $f \in  L_w^2(-1,1)$ Niech $N \in \mathbb{N}$\\
    Należy znaleźć współczynniki $\alpha_i \in \mathbb{R}$,  $i \in [N]$, \\
    takie, że, dla funkcji $f^* = \sum_{i=1}^{N}\alpha_i P_i$\\
    ~\\
    \centerline{$||f - f^*||$ jest najmniejsze.}\\
    ~\\
    Gdzie $P_i$ - wielomian Legendre'a $i-1$ stopnia\\

\end{alertblock}

\end{frame}


\begin{frame}

\frametitle{Wprowadzenie}
\begin{block}{Twierdzenie}

    Niech $V$ - przestrzeń Hilberta, $D<V$\\
    Element $f^* \in D$ jest elementem optymalnym dla $f \in V$\\
    względem podprzestrzeni D wtedy i tylko wtedy gdy\\
    \centerline{$\forall_{g \in D }$ <$f - f^*$,$g$> $= 0$.}
\end{block}

Twierdzenie to bardzo ułatwia rozwiązanie problemu.\\
Zakładająć, że $D = \mathbb{R}_{N-1}[x]$, oraz $D = \mathcal{L}(P_1, P_2, P_3, ... , P_{N-1}).$\\~\\
\centerline{Element $f^*$ jest optymalny wtedy i tylko wtedy gdy}
\centerline{$\forall_{g \in D }$ <$f - \sum\limits_{j=1}^{N}\alpha_j P_j$,$g$> $= 0$}\\
Jest to równoważne zdaniu:\\
\centerline{$\forall_{i \in [N-1] }$ <$f - \sum\limits_{j=1}^{N}\alpha_j P_j$,$P_i$> $= 0$}\\

\end{frame}

\begin{frame}
\frametitle{Wprowadzenie}

\centerline{<$f - \sum\limits_{j=1}^{N}\alpha_j P_j$,$P_i$> $= 0$}\\~\\
\centerline{$\iff$}\\~\\
\centerline{<$f$,$P_i$>  $-$  <$\sum\limits_{j=1}^{N}\alpha_j P_j$,$P_i$> $= 0$}\\
\centerline{$\iff$}\\~\\
\centerline{<$f$,$P_i$>  $= $<$\sum\limits_{j=1}^{N}\alpha_j P_j$,$P_i$>}\\~\\

Problem został sprowadzony do rozwiązania układu równań liniowych.\\

\end{frame}

\begin{frame}
\frametitle{Wprowadzenie}
\begin{block}{Twierdzenie}
$\forall_{i,j \in \mathbb{N}}$ $P_i * P_J = \frac{2}{2i-1}\delta_{ij}$

\end{block}
Powyższe twierdzenie znacząco upraszcza wyprowadzony układ równań.\\
\centerline{<$f$,$P_i$>  $= $<$\sum\limits_{j=1}^{N}\alpha_j P_j$,$P_i$>}\\~\\
\centerline{$\iff$}\\~\\
\centerline{<$f$,$P_i$>  $= \alpha_i * \frac{2}{2i-1} $}\\

\end{frame}

% Section 2
\section{Main Content}
\begin{frame}
\frametitle{Key Concepts}
\begin{block}{Definition}
    A block environment for definitions
\end{block}
\begin{alertblock}{Important}
    An alert block for warnings
\end{alertblock}
\begin{exampleblock}{Example}
    An example block
\end{exampleblock}
\end{frame}

\end{document}